% ============================================================
% === START OF PAPER: arXiv-2603.15759v1 ===
% === TITLE: Simulation Distillation: Pretraining World Models in Simulation for Rapid Real-World Adaptation ===
% ============================================================



\title{Simulation Distillation: Pretraining World Models in Simulation for Rapid Real-World Adaptation
\vspace{-0.25em}
}

\author{
  Jacob Levy$^{\ast1}$\quad Tyler Westenbroek$^{\ast2}$\quad Kevin Huang$^{2}$\quad Fernando Palafox$^{1}$\quad
  Patrick Yin$^2$\quad\\
  Shayegan Omidshafiei$^3$\quad Dong-Ki Kim$^3$\quad Abhishek Gupta$^2$\quad David Fridovich-Keil$^1$\\
  $^1$University of Texas at Austin\quad $^2$University of Washington\quad $^3$FieldAI \\
  \texttt{\textcolor{magenta}{\url{https://sim-dist.github.io}}}
  \thanks{
  $^\ast$Equal contribution.
  }
\vspace{-1.75em}
}



\maketitle


\begingroup
\setlength\stripsep{0pt}
\begin{strip}
    \centering
    \includegraphics[width=.99\textwidth]{figs/front.pdf}
    \captionof{figure}{\footnotesize{ Failures of zero-shot sim-to-real policies  (left). Our framework \Method\ rapidly overcomes the dynamics gap and improves performance with minimal real-world interaction. We demonstrate substantial gains in task execution on both precise manipulation and quadrupedal locomotion tasks with only 15-30 minutes of real-world data, substantially outperforming baselines. 
    }}
    \label{fig:front}
    \vspace{1em}
\end{strip}
\endgroup

\begin{abstract}
Simulation-to-real transfer remains a central challenge in robotics, as mismatches between simulated and real-world dynamics often lead to failures. While reinforcement learning offers a principled mechanism for adaptation, existing sim-to-real finetuning methods struggle with exploration and long-horizon credit assignment in the low-data regimes typical of real-world robotics. We introduce \texttt{Simulation Distillation} (\Method), a sim-to-real framework that distills structural priors from a simulator into a latent world model and enables rapid real-world adaptation via online planning and supervised dynamics finetuning. By transferring reward and value models directly from simulation, \Method\ provides dense planning signals from raw perception without requiring value learning during deployment. As a result, real-world adaptation reduces to short-horizon system identification, avoiding long-horizon credit assignment and enabling fast, stable improvement. Across precise manipulation and quadruped locomotion tasks, \Method\ substantially outperforms prior methods in data efficiency, stability, and final performance. Project website and code: \texttt{\textcolor{blue}{\url{https://sim-dist.github.io}}}
\end{abstract}

\IEEEpeerreviewmaketitle

\begin{figure*}[t]
    \centering
    \includegraphics[width=.99\textwidth]{figs/method.pdf}
    \caption{
    \footnotesize{\Method \ overview.
    1) An expert policy, policy checkpoints, and a value function are trained in simulation using privileged state. 2) Large-scale training data are generated by combining expert and sub-optimal policies with contiguous action perturbations, yielding diverse trajectories with dense reward and value supervision. 3) A planning-oriented latent world model is pretrained on this data, learning representations, dynamics, rewards, and values from raw observations. 4a) At deployment, the learned representation and dense reward and value models are transferred to the real robot to enable planning with the latent dynamics. 4b) Real-world data is then used to finetune only the dynamics via supervised system identification, with representations, rewards, and values frozen. Deployment and finetuning are iterated, enabling rapid and stable real-world adaptation.
    }
    }
     \vspace{-1em}
    \label{fig:method}
\end{figure*}

\section{Introduction}

Robotic systems must operate in complex, novel environments where success requires rapid deployment-time adaptation. Learning-based approaches \cite{intelligence2025pi,JMLR:v17:15-522, kim2024openvla} have shown great promise in this setting, but training robots directly in the real world remains challenging due to safety, cost, and time constraints.  As a result, simulators have emerged as a powerful tool for scaling robot learning, providing access to large volumes of interaction data and broad coverage over initial conditions and environment parameters \cite{rudin2022learning, yin2026emergent, narang2022factory}. While this strategy has enabled many recent success stories \cite{lee2020learning, kumar2021rma, singh2024dextrah}, sim-to-real failures remain common due to the inevitable dynamics ``gap'' between the two domains.

Bootstrapping behavior in simulation and finetuning in the real world with limited on-domain data to overcome dynamics mismatch is an attractive paradigm for robotics \cite{smith2022legged, yin2025rapidly,westenbroek2022lyapunovdesignrobustefficient,wagenmaker2024overcoming}. However, using existing end-to-end policy optimization approaches for this purpose often suffers performance collapse when finetuning in new domains \cite{zhou2024efficient, yin2025rapidly} due to catastrophic forgetting of the priors learned during pretraining. This reflects the difficulty of exploration and credit assignment in low-data regimes, particularly when learning from raw perception. Recent off-policy RL methods attempt to improve sample efficiency by aggressively bootstrapping from limited datasets \cite{rlpd, redq, droq, serl}, but face brittle trade-offs between data efficiency and training stability. We argue that this behavior stems from a fundamental limitation of end-to-end model-free policy finetuning: by tightly entangling task representations, rewards, and dynamics, these methods force algorithms to relearn the full task structure when transferred to new domains. In this work we ask: \emph{how can we retain physical priors and known task structure throughout finetuning?} 


We argue that world models, rather than end-to-end policies, are the appropriate abstraction for transferring experience from simulation to the real world. Unlike policies, world models decompose decision making into state representations, reward and value functions, and environment-specific dynamics that are learned jointly but maintained as distinct components. This decomposition enables an approximate separation between \emph{global} and \emph{local} task structure. State encoders, rewards, and values capture global structure—such as which objects matter, their semantic roles, and distances to goal configurations—while the dynamics model captures local, action-dependent behavior specific to the environment. Our key insight is that this global structure is largely invariant across simulation and reality for dynamic, high-precision problems, whereas local dynamics are not. We exploit this asymmetry to focus real-world adaptation on dynamics while preserving global task structure learned in simulation. 

Concretely, we introduce \texttt{Simulation Distillation} (\Method), a framework that bootstraps a world model in simulation and adapts in the real world through online planning and supervised finetuning of the latent dynamics model. At deployment, we freeze the state encoder, reward model, and value function learned in simulation, and adapt only the dynamics. The planner uses the frozen reward and value models to estimate  long-horizon returns, enabling immediate behavioral improvement as dynamics predictions become more accurate, while avoiding difficult bootstrapping problems. By planning over the world model, the robot can reason counterfactually about trajectories it has not directly experienced, mitigating exploration challenges. Altogether, this sidesteps the instability of existing off-policy methods by reducing real-world adaptation to a simple, short-horizon supervised learning problem in latent space. 

Enabling this reliable and efficient adaptation requires systematically distilling the structure of the simulator into the world model during pretraining. In particular, the dynamics model must generalize from only a small number of real-world samples and remain accurate along the many sub-optimal trajectories sampled by the planner. To achieve this, we build on standard teacher–student distillation frameworks, leveraging privileged, state-based expert policies in simulation to generate high-quality behavior at scale. Importantly, rather than just collecting expert data, we deliberately perturb actions during data generation. This introduces failures and recovery behaviors, yielding broad state–action coverage that keeps representations, rewards, values, and dynamics in-distribution during real-world planning and finetuning. Relying on privileged experts for pretraining enables massive parallelism and scalable training, in contrast to computationally intensive online model-based RL methods \cite{hafnerdream,hansen2022temporal} whose performance often saturates on challenging robotics problems.


We summarize our contributions as follows:


\begin{hangingpar}
\textbf{1)} We introduce \Method, a world-model–based paradigm for sim-to-real transfer that reduces real-world adaptation to a simple supervised dynamics learning problem, avoiding long-horizon credit assignment and enabling reliable, off-policy, and sample-efficient adaptation. \Method\ integrates seamlessly with existing student–teacher distillation.
\end{hangingpar}
\begin{hangingpar}
\textbf{2)} We instantiate \Method\ on two precise, contact-rich manipulation tasks and two quadruped locomotion tasks in the real world, achieving effective autonomous improvement with just 15–30 minutes of real-world data.
\end{hangingpar}
\begin{hangingpar}
\textbf{3)} We show that informative reward and value models trained in simulation can transfer zero-shot to the real world from raw perception, highlighting a promising direction for leveraging privileged simulation supervision even beyond the settings studied in this work.
\end{hangingpar}

\section{Related Work}

\textbf{Real-World Reinforcement Learning.} A growing body of work studies reinforcement learning on real-world robotic systems \cite{serl, levy2025learning, yin2025rapidly, smith2022walk, westenbroek2022lyapunovdesignrobustefficient, westenbroek2023enabling, JMLR:v17:15-522, pmlr-v87-kalashnikov18a, haarnoja2018soft, lanier2025adaptingworldmodelslatentstate}. However, both model-free and model-based approaches remain challenging to apply reliably in low-data regimes and typically require sophisticated regularization. Model-free methods aggressively reuse off-policy data and rely on frequent critic updates \cite{rlpd, droq, serl, redq}, often leading to value overestimation and unstable learning. Prior work mitigates these issues using conservative value estimation \cite{kumar2020conservativeqlearningofflinereinforcement}, policy constraints \cite{lei2025rl, zhang2025rewindlanguageguidedrewardsteach, smith2024grow}, or critic ensembles \cite{chen2021randomized, haarnoja2018learning, hiraoka2022dropout}. Model-based methods instead learn dynamics and reward models \cite{hafnerdream, hansen2022temporal} to reason about unseen trajectories, but must carefully avoid exploiting model inaccuracies. Common strategies include uncertainty-aware dynamics \cite{chua2018deep, levy2025learning, janner2021trustmodelmodelbasedpolicy} and penalizing out-of-distribution predictions \cite{yu2020mopo, yu2021combo, feng2023finetuningofflineworldmodels}. Rather than constraining learning, we show that bootstrapping a world model in simulation enables generalization beyond the real-world dataset.


\textbf{Adaptation with Physics-based Models.} Many lines of work leverage approximate physics-based models for adaptation and control, but typically rely on simplified, low-dimensional state representations that are brittle in partially observed, contact-rich settings. Classical adaptive control \cite{aastrom1995adaptive, slotine1987adaptive, ioannou1996robust} and model-predictive control \cite{morari1999model, 8594448} approaches use highly simplified dynamics, abstracting away complex contacts and interactions. Neural physics engines \cite{xu2025neural, pfrommer2021contactnets} combine structured system identification with residual learning to adapt high-fidelity simulators using real-world data, but often assume access to object poses, contact labels, or reliable state estimates that degrade under partial observability. Closely related work learns world models from mixtures of simulated and real data \cite{li2025offline, xiao2025anycar} or transfers value functions from simulation to guide real-world learning \cite{westenbroek2022lyapunovdesignrobustefficient, yin2025rapidly}, but similarly depends on low-dimensional state observations. We instead propose distilling simulator structure from raw perception to enable in-the-wild adaptation. 



\textbf{Generative World Models for Robotics.} Recent work trains large  world models on internet-scale video to learn broad physical priors for robotics \cite{yang2023learning, bruce2024genie, parkerholder2024genie2}, sometimes augmented with simulation data \cite{yang2023learning}. Translating video predictions into executable robot actions typically requires expert demonstrations, either by planning in video space and using inverse models to recover actions \cite{xie2025latent, jang2025dreamgen}, or by combining video prediction with behavior cloning \cite{zhu2025unified, li2025unified, cheang2024gr}. While effective for reproducing demonstrated behaviors, these approaches remain fundamentally constrained by the real-world action distribution. In contrast, \Method\ learns world models directly over a wide range low-level robot actions, enabling it to directly improve behaviors beyond the real-world data.

\section{ Preliminaries}

\textbf{Problem Setting.} We consider the problem of controlling a robotic system operating in the real world under partial observability and unmodeled dynamics. Specifically, we assume the dynamics are of the form $s_{t+1} \sim p(\cdot | s_t,a_t)$ where $s_t \in \mathcal{S}$ is the underlying state of the system and $a_t \in \mathcal{A}$ is the robot action. The underlying state is not directly observable in the real world as the robot only has access to raw observations $o_t \in \mathcal{O}$. We model the control problem as a partially observable Markov decision process (POMDP) defined by the tuple $(\mathcal{S}, \mathcal{A}, \mathcal{O}, \mathcal{X}, p, r, \gamma)$, with user-specified reward function $r(s_t, a_t, s_{t+1})$ and discount factor $\gamma \in \left(0, 1\right]$, and seek to maximize the discounted return $\mathbb{E} \left[ \sum_{t=0}^{\infty} \gamma^t r(s_t, a_t, s_{t+1})\right]$. In general, dense informative reward functions are difficult to evaluate directly in the real world, given the difficulty in measuring the underlying state of the system. To render this problem more tractable, we make the following assumption: 
\begin{assumption}
We assume access to an approximate physics-based simulator $s_{t+1} \sim \hat{p}(\cdot \mid s_t, a_t)$ that provides privileged access to the underlying state $s_t$. 
\end{assumption}

\begin{figure}[t]
    \centering
    \includegraphics[width=.48\textwidth]{figs/simdist_architecture.pdf}
    \caption{
    \footnotesize{
    World model architecture.
    The most recent observation is encoded into a latent representation while a history encoder processes a history of observations and actions. These jointly condition a transformer-based latent dynamics model that predicts future latent trajectories under candidate action sequences. Transformer-based reward and value heads evaluate predicted trajectories to produce reward and value sequences, while a base policy head predicts action chunks used to warm-start sampling-based planning.}
    }
    \vspace{-1em}
    \label{fig:model}
\end{figure}


\textbf{Planning-Oriented Latent World Models.} We build on common planning-oriented latent world models from works such as \cite{hansen2022temporal, hafnerdream}. These approaches learn reward and value models that operate directly on raw observations, enabling planning in the real world using the algorithms outlined below. The primary novelty of \Method \ lies in our systematic framework for sim-to-real adaptation of this class of world models, introduced in the following section. We structure our world model according to \cref{eq:model}, depicted in \cref{fig:model}:
\begin{equation}
\label{eq:model}
\begin{array}{ll}
    \text{Latent representation:} & \latent_{t} = \encoder_{\theta}(o_t)\\
    \text{History representation:} & h_t = \contextencoder_{\theta}(o_{t-H:t-1}, a_{t-H:t-1})\\
    \text{Latent dynamics:} & \hat{\latent}_{t+1:t+T} = \dynamicsmodel_{\theta}(\latent_{t}, \action_{t:t+T -1}, h_t)\\
    \text{Reward Prediction:} & \hat{\reward}_{t:t+T-1} = \rewardmodel_{\theta}(\hat{z}_{t:t+T},a_{t:t+T-1})\\
    \text{Value Prediction:} & \hat{\val}_{t+1:t+T} = \valuemodel_{\theta}(\hat{\latent}_{t:t+T})\\
    \text{Base Policy:} & \hat{\action}_{t:t+H} = \bcpolicy_{\theta}(\latent_{t}, h_t).
\end{array}
\end{equation}
Here $E_{\theta}$ encodes a latent representation $z_t$ of the underlying state $s_t$ directly from raw observation $o_t$; $h_t$ is an encoding of history over a window of $H$ timesteps; $\hat{z}_{t+1:t+T}$ is a predicted sequence of $T$ future latent states generated by the learned model $f_\theta$; and $\hat{r}_{t:t+T-1}$ and $\hat{v}_{t+1:t+T}$ are reward and value estimates. We discuss key architecture decisions in \cref{sec:modeldetails}. 

\textbf{Sampling-Based Planning.}
We control the robot in the real world with \ac{mppi} \cite{williams2016aggressive} control, a sampling-based \ac{mpc} method. At each control step, we sample a batch of candidate future action sequences and evaluate them using the world model \cref{eq:model}, extracting the predicted future reward sequence $\hat{r}_{t:t+T-1}$ and terminal value $\hat{v}_{t+T}$, which are combined to compute the trajectory return \mbox{$\mathcal{R}(a_{t:t+T-1}) = \gamma^T\hat{v}_{t+T} + \sum_{s=t}^{t+T-1} \gamma^{s-t}\hat{r}_s$}.
MPPI then computes the control action to execute by importance-weighting sampled trajectories according to their predicted returns. Similar to TD-MPC \cite{hansen2022temporal}, we warm-start sampling by seeding a subset of candidate action sequences with noise-corrupted outputs of $\hat{a}_{t:t+T-1}$ from the base policy $\pi_\theta$.


\section{Simulation Distillation for Efficient Real-World Adaptation}

We now introduce \texttt{Simulation Distillation}, a framework for distilling physical priors and task structure into a world model in a form that enables rapid real-world adaptation with online planning and dynamics adaptation.  


\subsection{Pretraining on Simulated Data} \label{sec:pretrain}
 
\Method\ builds on existing sim-to-real data-generation pipelines that use privileged, state-based expert policies to collect large-scale datasets paired with raw perception. While prior approaches primarily use this data for imitation, planning-based adaptation requires world models that can reason counterfactually about trajectories far beyond the expert and real-world data distributions. In particular, dynamics must generalize from few real-world samples, and reward and value models must reliably discriminate between good and bad trajectories under off-policy action sequences. To meet these requirements, we deliberately inject sub-optimal actions during simulation rollouts, producing a broad training distribution.

\textbf{Expert Policy Training.} We first train a state-based expert policy $\pi^e(s_t)$ with PPO \cite{schulman2017proximal}. We also save the optimal state-based value function $V^e(s_t)$ and intermediate policy checkpoints $\{\pi^k\}_{k=1}^{K}$ learned during training.


\textbf{Generating Diverse Trajectories and Dense Supervision.} We generate diverse simulation rollouts (Fig.~\ref{fig:method}) by alternating between an expert policy $\pi^e$ and a set of sub-optimal policies $\{\pi^k\}_{k=1}^K$, and by periodically injecting random action perturbations over short temporal windows. This generates diverse failure and recovery behaviors beyond the nominal expert manifold, enabling the world model to reliably produce predictions about counterfactual trajectories. This results in a dataset $\simdataset = \{(o_t, a_t, r_t, v_t)\}_{t=0}^N$, where rewards $r_t$ are computed from privileged simulator state and value targets $v_t$ are provided by the expert value function $V^e(s_t)$. Crucially, this data-generation process is massively parallelizable and exploits the full training artifact of the simulator—including expert policies, intermediate checkpoints, and value functions—to produce rich supervision at scale. The full data generation procedure is provided in \cref{app:datagendetails}.



\textbf{World Model Pretraining.}
We pretrain the world model \cref{eq:model} on $\mathcal{D}_{\text{sim}}$ by applying the following loss to predictions made at each time step $t$. Specifically, we feed the history $o_{t-H:t}$ and future actions $a_{t:t+T-1}$ into the model, and then penalize the resulting predictions with the following loss: 
\begin{equation}
\label{eq:simloss}
\begin{aligned}
\vspace{-.1em}
    \mathcal{L}_t^{\mathtt{sim}}(\theta) = \sum_{i=0}^{T}\bigg(\,\, &\underbrace{\norm{\hat{z}_{t+i+1} - \mathtt{sg}(E_{\theta}(o_{t+i+1}))}^2_2}_{\text{latent dynamics}}\\
    &+ c_1\underbrace{(\hat{r}_{t+i} - r_{t+i})^2}_{\text{reward}} + \, c_2\underbrace{(\hat{v}_{t+i+1} - v_{t+i+1})^2}_{\text{value}}\\
    &+ c_3 \underbrace{\mathds{1}_e(a_{t+i}) \norm{\hat{a}_{t+i} - a_{t+i}}^2_2}_{\text{behavior cloning}}\bigg),
\end{aligned}
\end{equation}
where $\mathtt{sg}$ is the \texttt{stop-grad} operator, constants $c_{1:3}$ are determined by normalizing over the range for each target to ensure each loss term is roughly on the same scale, and $\mathds{1}_e(a_t)=1$ if $\action_t$ came from the uncorrupted expert policy and $\mathds{1}_e(a_t)=0$ otherwise. We apply various data augmentations discussed in \cref{app:manip,app:qped} to prevent overfitting to simulated observations an ensure that the model learns robust, transferable representations. In contrast to typical online MBRL approaches \cite{hansen2022temporal,hafnerdream} which are computationally intensive and attempt to bootstrap behavior from scratch, we offload behavior generation to a privileged expert. This reduces pretraining to optimizing a simple, stationary objective over a fixed, large-scale, and diverse simulated dataset.

\subsection{Real World Transfer and Efficient Dynamics Adaptation}
\label{sec:finetune}
Our key insight is that global task structure is largely invariant to low-level sim-to-real dynamics gaps. For example, in Peg Insertion (\cref{fig:front}), a meaningful latent state captures the locations of the peg and hole, while the value function encodes distance to the goal and motions leading to successful insertion. This structure persists across the sim-to-real gap, even though the low-level actions required to realize these behaviors differ between domains. Exploiting this decomposition, \Method\ finetunes only the environment-specific dynamics model while freezing the encoder, reward, and value functions. Planning with frozen reward and value models enables immediate improvement as dynamics predictions become more accurate, without requiring reward or value bootstrapping in the real world. Because the world model can be trained on arbitrary trajectories, \Method\ naturally supports off-policy learning and can incorporate prior data such as demonstrations.


\textbf{Dynamics Adaptation Loss.} When updating the dynamics model, we modify the pretraining loss for predictions made at each time-step $t$ to simply: 
\begin{equation}
\label{eq:realloss}
\begin{array}{l}
    \mathcal{L}_t^{\mathtt{real}}(\theta) = \sum_{i=0}^{T} \,\, \norm{\hat{z}_{t+i+1} - \mathtt{sg}(E_{\theta}(o_{t+i+1}))}^2_2, \\
    \text{with } C_\theta, E_\theta, R_\theta, V_\theta, \pi_\theta \ \text{frozen, } f_{\theta} \text{ finetunable}. 
\end{array}
\end{equation}
We do not apply data augmentations when optimizing this objective, as pretraining yields a robust encoder $E_\theta$. Because the encoder is frozen, it provides consistent latent targets 
$E_\theta(o_t)$ throughout finetuning, avoiding the need to boot strap a representation as in \eqref{eq:simloss}. 



\textbf{Iterative Improvement.} The overall pipeline for \Method \ is shown in \cref{fig:method} and in pseudo-code in \cref{alg:method}. \Method \ autonomously improves in the real-world by repeatedly collecting $M$ on-policy rollouts under the planner, adding this data to the real-world data set $\realdataset$, then finetuning $f_\theta$ to minimize prediction losses as in \eqref{eq:realloss}. Notably, because system identification can effectively learn from any real-world trajectories, \Method \ is a simple off-policy reinforcement learning strategy which can easily incorporate diverse data sources such as demonstrations or play data into $\realdataset$.

\begin{remark}
In contrast to standard MBRL frameworks \cite{hafnerdream,hansen2022temporal}, which must jointly bootstrap latent representations, value functions, and policies from scarce in-domain data, \Method\ offloads these challenging objectives to simulation, where diverse data is cheap and plentiful. As a result, real-world adaptation reduces to finetuning only the dynamics via a simple, iterative supervised learning procedure.
\end{remark}
\subsection{World Model Design Decisions}
\label{sec:modeldetails}

In order to successfully improve behavior over the base policy $\pi_\theta$, the planner must sample numerous off-policy rollouts and accurately model returns. The following design decisions are essential for enabling reliable real-time planning. 

\textbf{Minimal History Representation.}
Observation histories $o_{t-H:t}$ can contain high-dimensional inputs such as images that are costly to process. To reduce inference cost, we split observations $o_{t} = (o_t^p,o_t^e)$ into proprioceptive and exteroceptive components and feed only $(o_{t-H:t}^p,a_{t-H:t},o_t^e)$ into the history encoder $C_\theta$. Using only the most recent high-dimensional observation substantially reduces planning latency and, empirically, improves training stability by reducing context length.

\textbf{Chunked Prediction for Planning.} Autoregressive world models \cite{hansen2022temporal, hafnerdream, chen2022transdreamer} require sequential unrolling over the planning horizon, which bottlenecks parallelism when  evaluating numerous rollouts. Inspired by \cite{xiao2025anycar}, our latent dynamics model $f_\theta$ predicts $T$ future states in a single forward pass using a transformer with cross-attention between the encoded history tokens and a candidate action sequence, together with a causal mask. This chunked prediction fully exploits GPU parallelism, dramatically improving planning throughput. 

\textbf{Sequence-to-Sequence Return Modeling.} Prior work often decodes rewards and values with per-timestep MLPs applied independently to each predicted latent state \cite{hansen2022temporal, hafnerdream, chen2022transdreamer}. We instead use transformer-based reward and value models that attend over the entire predicted latent trajectory $\hat{z}_{t:t+T}$. This increased capacity aggregates information across the entire trajectory, yielding more accurate return estimates and improving planning performance (see \cref{sec:ablations}).


\section{Experiments}
\label{sec:exp}

\begin{figure*}
    \centering
    \includegraphics[width=0.99\textwidth]{figs/robots_for_results.pdf}
    \includegraphics[width=0.99\textwidth]{figs/results.pdf}
    \vspace{-1em}
    \caption{\footnotesize{
    Real-world results. Success rate for two manipulation tasks, computed over 20 trials, and average forward progress for two quadruped locomotion tasks, averaged across all $15$ trials (3 speeds, 5 trials each), as a function of real-world finetuning data. \Method\ exhibits rapid and consistent improvement with limited data by finetuning only the latent dynamics model while planning with frozen reward and value models. In contrast, direct policy finetuning with the baselines shows limited or no improvement under the same data budgets.
    }}
    \vspace{-1em}
    \label{fig:results}
\end{figure*}

We evaluate the ability of \Method\ to adapt in the real world on the four manipulation and quadruped tasks depicted in \cref{fig:front} and carefully ablate key design decisions in simulation. Additional details of experiments can be found in the Appendix. We aim to answer $1)$ does \Method \ outperform existing online RL methods and behavior cloning baselines? $2)$ what factors enable the planner to effectively improve performance with minimal real-world data?  $3)$ which components of the architecture and pretraining procedure are crucial for the performance of \Method?  


\subsection{Robotic Systems and Tasks}
\textbf{Manipulation System and Tasks. } Our manipulation experiments are conducted on a UR5e robot. Actions are six-dimensional relative end-effector pose targets, and observations include joint states and three $224\times224$ RGB images from wrist-mounted, overhead, and side-view cameras. Each image is encoded with a pretrained ResNet-18, fused with proprioception, and mapped to a 64-dimensional latent state $z$. Training uses 100K trajectories, approximately $36\%$ of which are optimal expert sequences (see \cref{app:manip} for details). We use history and prediction horizons $H=T=5$ and control the robot at \SI{5}{\hertz}. Expert policies are trained following \cite{yin2026emergent}. We consider the following precise assembly tasks, with initial conditions drawn from a Narrow (\qtyproduct{2 x 2}{\cm}) or Wide (\qtyproduct{35 x 35}{\cm}) grid:
\begin{hangingpar}
\textbf{1) Peg Insertion.} We construct a peg insertion task similar to the \SI{16}{\mm} square peg task from \cite{narang2022factory}. This task requires picking the peg, aligning it with the hole, and insertion.
\end{hangingpar}
\begin{hangingpar}
\textbf{2) Table Leg.} We use the assets from \cite{heo2023furniturebench} to generate a task where a table leg must be picked, aligned, and threaded into a hole on a table.
\end{hangingpar}
\textbf{Quadruped System and Tasks.} We conduct quadrupedal experiments on a Unitree Go2 platform, with actions given as position targets for the $12$ joints. Observations include proprioceptive signals and a local terrain height map, which is encoded by $E_\theta$ using a CNN and fused with low-dimensional observations via an MLP to produce the latent state $z$. The policy $\pi_\theta$, reward $R_\theta$, and value $V_\theta$ are additionally conditioned on desired base forward, lateral, and yaw velocities. Pretraining uses 100M simulated trajectories, of which $55.7\%$ come from an expert policy trained in IsaacLab \cite{mittal2025isaaclab}. We use history and prediction horizons $H=T=25$ and plan at \SI{50}{\hertz} on a laptop with an RTX 4090M. See \cref{app:qped} for details. We consider the following tasks: 
\begin{hangingpar}
\textbf{1) Slippery Slope.} The robot traverses straight over two panels inclined at $3.0^\circ$ and $5.7^\circ$, respectively. The panels are covered with PTFE (Teflon), and the robot’s feet are wrapped with thermoplastic, creating extremely low-friction contacts. A trial is successful if the robot traverses \SI{1.82}{\meter}, clearing both panels. We conduct five trials per commanded forward speed at $0.1$, $0.3$, and $0.5$~\si{\meter\per\second}.
\end{hangingpar}

\begin{hangingpar}
\textbf{2) Foam.} The robot traverses two \SI{5}{\cm} thick overlapping memory foam pads whose compliant dynamics are not modeled in simulation. A trial is  successful if the robot traverses \SI{3.00}{\meter} to clear the foam. We conduct five trials per commanded forward speed at $0.2$, $0.7$, and $1.2$~\si{\meter\per\second}.
\end{hangingpar}

\subsection{Real-World Learning Methods.}

We evaluate \Method\ against state-of-the-art real-world RL methods and additionally behavior cloning baselines on manipulation tasks. All RL methods use the same encoders as \Method, but with an MLP policy head (no action chunking). 

\textbf{Manipulation.}
For each task, we collect $20$ real-world demonstrations via teleoperation. We evaluate two variants of \Method: (i) \Method, which adapts only the latent dynamics model using on-policy rollouts (\cref{alg:method}), and (ii) \Method+BC, which additionally finetunes the base policy $\pi_\theta$ using expert action labels. In both cases, the dynamics model is updated after every 20 real-world episodes. We compare against model-free RL baselines RLPD \cite{rlpd} and IQL \cite{kostrikovoffline}, trained from sparse rewards to reflect common manipulation settings where dense rewards are difficult to specify; these baselines update after every episode. We also include SGFT-SAC \cite{yin2025rapidly}, a model-free method that transfers a simulated value function, isolating the benefit of full world-model adaptation beyond value transfer alone. Finally, we compare against behavior cloning baselines—Diffusion Policy \cite{chi2024diffusionpolicy} and $\pi_{0.5}$ \cite{intelligence2025pi05}—trained on real-world demonstrations, as well as variants trained on both real-world demonstrations and the simulated dataset used by \Method. All manipulation methods except \Method\ without demonstrations are given access to the same $20$ real-world demonstrations.

\textbf{Quadruped.}
For quadruped locomotion, we compare \Method\ against the off-policy algorithm RLPD \cite{rlpd} and the offline-to-online finetuning method from IQL \cite{kostrikovoffline}, whose value function is pretrained in simulation. All methods use the same learned reward model from \Method, allowing us to isolate the effect of different adaptation strategies.

\subsection{Real World Improvement Results}

\Cref{fig:results} summarizes our results. Across all tasks, \Method\ consistently outperforms prior approaches, achieving substantially higher success rates with far greater sample efficiency than online RL baselines, while autonomously improving well beyond the performance of behavior cloning methods. Across the board \Method \ typically reaches scores around $2\times$ higher than any baseline. Standard RL finetuning methods frequently exhibit catastrophic forgetting, with performance collapsing during adaptation, whereas \Method\ makes steady, monotonic progress throughout training. SGFT avoids catastrophic collapse by transferring value functions from simulation, but remains significantly more sample inefficient than \Method, highlighting the importance of adapting a full world model. Interestingly, demonstrations provide a larger boost to \Method \ on the Leg task, where the precise screwing motion was difficult to transfer. This highlights how \Method \ can naturally leverage demonstrations when available. 

Notably, the performance gap between \Method\ and all baselines widens on more challenging tasks, particularly when comparing narrow versus wide ranges of initial conditions. This trend underscores the benefit of broad simulation pretraining, which enables \Method\ to generalize across diverse initial states and adapt reliably in difficult real-world regimes. This effect is illustrated in \cref{fig:scatter}, which visualizes successful and failed initial conditions for \Method\ and Diffusion Policy on Peg Hard, revealing the substantially greater robustness of policies learned by \Method.

\begin{figure}
    \centering
    \includegraphics[width=0.48\textwidth]{figs/value_overlay.pdf}
    \vspace{-0.5em}
    \caption{\footnotesize{Value predictions from \Method\ along successful and failed real-world Peg trajectories starting from the same initial condition. The predicted values track task progress and clearly distinguish successful from failure.}}
    \label{fig:value}
\end{figure}


\begin{figure}
    \centering
    \includegraphics[width=0.48\textwidth]{figs/scatter.png}
    \vspace{-0.5em}
    \caption{
    \footnotesize{Scatter-Plot showing successful and failed attempts at solving the Peg Wide task for Diffusion Policy (Right) and the final trained policy for \Method (Left). The broad coverage of pretraining data for \Method \ enables efficiently learning policies which are far more robust than baselines.}
    }
    \label{fig:scatter}
    \vspace{-1em}
\end{figure}

\begin{figure*}
     \centering
     \begin{subfigure}[t]{0.30\textwidth}
         \centering
         \panel{figs/consistency_loss_plot.pdf}{(a)}{\linewidth}
         \phantomsubcaption
         \label{fig:consistloss}
     \end{subfigure}
     \hfill
     \begin{subfigure}[t]{0.140\textwidth}
         \centering
         \panel{figs/foot_slip.pdf}{(b)}{\linewidth}
         \phantomsubcaption
         \label{fig:footslip}
     \end{subfigure}
     \hfill
     \begin{subfigure}[t]{0.235\textwidth}
         \centering
         \panel{figs/foot_preds.pdf}{(c)}{\linewidth}
         \phantomsubcaption
         \label{fig:footpreds}
     \end{subfigure}
     \hfill
     \begin{subfigure}[t]{0.295\textwidth}
         \centering
         \panel{figs/planning.pdf}{(d)}{\linewidth}
         \phantomsubcaption
         \label{fig:planning}
     \end{subfigure}
     \hfill
     \vspace{-1em}
     \caption{\footnotesize{
     \textbf{(a)} Finetuning drastically lowers dynamics prediction loss during a quadruped Slippery Slope trial.
    \textbf{(b)} Overlaid frames showing the front left foot slipping during the trial.
    \textbf{(c)} Foot-trajectory predictions from the world model at the same instant: the finetuned model correctly anticipates future slippage, while the pretrained model fails to do so.
    \textbf{(d)} Visualization of sampling-based planning. Candidate action sequences are evaluated with the world model to produce planned foot trajectories for each leg. Sampled trajectories are shown from the finetuned dynamics model, along with the resulting optimal plans under the finetuned and pretrained models. The finetuned model produces plans that account for real world dynamics mismatch, while the pretrained model generates qualitatively different trajectories.
    }}
    \vspace{-1em}
\end{figure*}


\subsection{Analyzing Real-World Predictions and Effects on Planning}
\label{sec:ablations} The planner used by \Method\ can only improve behavior if it can reliably distinguish trajectories with high and low returns. This requires both accurate dynamics predictions and successful transfer of reward and value models. We examine value transfer in \cref{fig:value}, which plots predicted values over time for successful and failed trajectories, both collected with teleoperation from the same initial condition. For the successful rollout, predicted value increases consistently over time, while remaining lower for the failed trajectory. Thus the transferred encoder $E_\theta$ and value function $V_\theta$ can \textbf{reliably discriminate between successful and failed trajectories}.

We next evaluate dynamics prediction accuracy on the Quadruped Slippery Slope task in \cref{fig:consistloss}. We roll a held-out real-world trajectory through the world model and compute the latent dynamics loss \cref{eq:realloss} at each timestep, yielding an average loss of $0.076$ for the pretrained model and $0.019$ after finetuning. To visualize the impact of this improvement, we decode predicted latent states into predicted front-left foot positions in \cref{fig:footpreds}. Before finetuning, the model incorrectly predicts stable contact on the PTFE surface and fails to anticipate slip; after finetuning, it accurately predicts future slippage, closely matching the real trajectory (\cref{fig:footslip}). This illustrates how broad simulation pretraining enables the dynamics model to generalize to real-world trajectories outside the training set. 

Finally, improved dynamics predictions directly reshape planning behavior. As shown in \cref{fig:planning}, trajectory samples generated with the finetuned model reflect the altered contact dynamics and lead the planner to select plans that account for real-world slip. In contrast, plans derived from the pretrained model are qualitatively inconsistent with the true dynamics. Together, these results show that finetuning corrects latent dynamics errors in a way that meaningfully changes the planner’s trajectory distribution, explaining the rapid real-world performance gains observed in \cref{fig:results}.

\subsection{Ablating Key Design Decisions}





We ablate key components of \Method\ by modifying world-model pretraining and evaluating the resulting closed-loop planning policies in simulation. Results are summarized in \cref{tab:ablations}, where we evaluate policies under a wide range domain randomization for both manipulation and quadrupeds (see \cref{app:ablations} for details). 

\textbf{Scaling Simulation Data.}
We study the effect of dataset scale by reducing simulation rollouts to $10\%$ and $50\%$ of the full pretraining data. Performance drops sharply as data is reduced for both systems, indicating that large-scale pretraining is critical for achieving the coverage required for reliable planning.

\textbf{Data Diversity.}
We next evaluate the role of data diversity by comparing pretraining on expert-only trajectories to our mixed datasets with equal data volume. Expert-only data leads to a substantial performance drop, highlighting the importance of diverse, perturbed rollouts for learning dynamics and value estimates that remain accurate under off-policy planning.

\textbf{Reward and Value Transformers.}
We ablate the use of sequence-to-sequence transformers for reward and value prediction, replacing them with per-timestep MLP decoders as in \cite{hafnerdream, hansen2022temporal}. This consistently degrades performance across tasks. We attribute this to the inability of per-step models to capture trajectory-level structure, which is essential for accurately evaluating candidate action sequences during planning.

\textbf{Reconstruction Loss.}
Finally, we study the effect of adding an observation reconstruction loss to the dynamics objective, a common design choice in MBRL \cite{hafnerdream, bruce2024genie}. While reconstruction slightly improves quadruped performance, it provides no qualitative benefit and leads to a steep performance drop for manipulation. We hypothesize that reconstructing high-dimensional visual observations encourages overfitting to task-irrelevant details, which interferes with learning compact latent dynamics suitable for planning.

\section{Conclusion}
We introduced \Method, a framework for sim-to-real transfer and real-world adaptation that exploits the modular structure of world models to directly target the dynamics gap between simulation and reality. Across two precise manipulation tasks and two quadruped locomotion tasks, \Method\ significantly outperforms prior approaches in both sample efficiency and robustness, achieving consistent improvement where standard RL finetuning methods often collapse. While our experiments focus on single-task settings with high-fidelity simulators, future work will scale \Method\ to multi-task world models and broader domains. More broadly, freezing representations and task structure while adapting only dynamics offers a general and effective paradigm for efficient adaptation beyond sim-to-real transfer.



\section*{Acknowledgments}
The authors would like to thank Trey Smith and Brian Coltin for their helpful insights and feedback. This work was supported by a NASA Space Technology Graduate Research Opportunity under award 80NSSC23K1192.

\bibliographystyle{plainnat}
\bibliography{references}

\appendix
\crefalias{subsection}{appendix}

\subsection{Diverse Data Generation Details}
\label{app:datagendetails}
\Cref{alg:datagen} details the data generation process. Data generation proceeds with running many environments in parallel. For the $j$-th environment, we sample a diagonal action-noise covariance $\Sigma_j$, where each element on the diagonal is sampled between a minimum and maximum variance: $\sigma_i \sim \mathcal{U}[\sigma^\mathtt{min}_i, \sigma^\mathtt{max}_i]$. When each environment is reset, we sample contiguous noise intervals during which Gaussian noise is added to policy actions. In addition, each environment is assigned a randomly selected policy checkpoint from $\{\pi^k\}_{k=1}^{K}$, or the expert policy $\pi^e$. Together, policy mixing and temporally structured action noise produce a dataset that captures both expert-like trajectories and systematic deviations beyond the optimal manifold, which is critical for  world-model training. During rollouts, we query the optimal state-based value function $V^e$ at each visited state to generate value targets $v_t$ for distilling an approximate optimal value function (\cref{sec:pretrain}). We also record an expert action flag $\expertflag_t$ that distinguishes the expert actions from the expert policy $\pi^e$ versus noised or earlier-checkpoint actions to support behavior cloning (\cref{sec:pretrain}).

\subsection{Manipulation Experiment Details}
\label{app:manip}
\textbf{Expert Policy Training.} For our manipulation experts we use the expert policies $\pi^e$ and value function $V^e$ from \cite{yin2026emergent}, replicating the training from this work exactly. For the sake of brevity, we refer the reader to \cite{yin2026emergent} for more details, including the exact rewards we use. 


\textbf{Data Generation.} To generate the simulation dataset $\simdataset$, environments are reset with sub-optimal policies $\{\pi^k\}_{k=1}^K$ with probability $0.5$, and Gaussian action perturbations are injected in contiguous intervals sampled from $\mathcal{U}[1, 5]$ steps, interleaved with noise-free intervals sampled from $\mathcal{U}[5, 10]$ steps. We generate 100k trajectories for each tasks, of which approximately 36\% optimal actions. We save policies every 100 checkpoints up to checkpoint 1000. 


\textbf{World Model Structure.} The encoder for the world model first passes each of the three camera images through a ResNet-18 encoder pretrained on imagenet, producing embeddings of size $3 
\times 512$, which are stacked and concatenated with the robots 6 joint observations and then passed through an MLP to produce the latent $z$. Specific parameters are in  \cref{tab:model_man}.


\begin{table}[htb]
\caption{\footnotesize{World model architectural parameters for manipulation.}}
\centering
\begin{tabular}{lc}
\hline
Parameter & Value \\
\hline
Embedding dimension & 64 \\
All transformers MLP hidden size & 256 \\
Dynamics transformer layers & 3 \\
Dynamics transformer heads & 4 \\
Reward transformer layers & 1 \\
Reward transformer heads & 1 \\
Value transformer layers & 1 \\
Value transformer heads & 1 \\
Base policy transformer layers & 4 \\
Base policy transformer heads & 8 \\
\hline
\end{tabular}
\vspace{0.05cm}
\label{tab:model_man}
\end{table}




\textbf{World Model Pretraining.}
We pretrain the world model for two epochs over the full simulation dataset $\simdataset$. Using a batch size of 256 and approximately 200k gradient updates. Optimization is performed with Adam using an initial learning rate of \num{2e-4}, which is annealed to \num{1e-4} via a cosine decay schedule, with a linear warmup over the first 10{,}000 steps. We apply data augmentation by injecting zero-mean Gaussian noise into the proprioceptive observations, and visual augmentations such as color jitter,  gaussian blurring and random cropping. 


\textbf{Hardware Deployment.} We use MPPI as implemented in TD-MPC \cite{hansen2022temporal} with the hyperparameters listed in \cref{tab:mppi_manip}. 


\begin{table}[htb]
\caption{\footnotesize{MPPI parameters for manipulation.}}
\centering
\begin{tabular}{lc}
\hline
Parameter & Value \\
\hline
Candidate actions batch size & 250 \\
Noised base policy actions batch size & 100 \\
Solver iterations & 3\\
Initial action standard deviation & 1.0 \\
Minimum action standard deviation & 0.05 \\
Base policy action standard deviation & 0.1 \\
Elites & 64 \\
Temperature & 0.4 \\
Momentum & 0.0 \\
Discount & 0.99 \\
\hline
\end{tabular}
\vspace{0.05cm}
\label{tab:mppi_manip}
\end{table}




\subsection{Quadruped Experiment Details}
\label{app:qped}

\textbf{Expert Policy Training.} We train a state-based expert policy $\pi^e$ and its associated optimal value function $V^e$ using PPO in IsaacLab \cite{mittal2025isaaclab}. Both the policy and value networks are MLPs with three hidden layers of width 512 and operate on privileged simulator state variables listed in \cref{tab:state_qped}. The expert is trained with a dense state-based reward composed of the terms summarized in \cref{tab:reward_qped}; full implementation details will be released with the public code. To improve robustness and coverage, we apply domain randomization over the parameters in \cref{tab:state_qped} and randomize terrain conditions across steps, boxes, rough terrain, and slopes, using a curriculum that gradually increases terrain difficulty. Training is performed with 4096 parallel simulation environments over 5000 PPO iterations, with 24 environment steps per iteration, for a total of 490M environment steps. We save policy checkpoints at iterations $\{0, 50, 100, 150, 200, 250, 300, 400, 500, 1000, 2000\}$.


\textbf{Data Generation.} To generate the simulation dataset $\simdataset$, environments are reset with sub-optimal policies $\{\pi^k\}_{k=1}^K$ with probability $0.5$, and Gaussian action perturbations are injected in contiguous intervals sampled from $\mathcal{U}[1, 50]$ steps, interleaved with noise-free intervals sampled from $\mathcal{U}[25, 500]$ steps. We run 4096 parallel environments for 25000 steps, yielding approximately 100M data points, of which $55.7\%$ correspond to uncorrupted expert actions. Data generation takes approximately 7 hours with a single NVIDIA RTX 4500 Ada GPU.

\textbf{World Model Structure.} \Cref{fig:model_qped} illustrates the world model architecture, and \cref{tab:model_qped} lists the corresponding model parameters used in the quadruped experiments. The observation space for the quadruped is given in \cref{tab:obs_qped}. The history encoder processes a history of proprioceptive observations (all observations except the height map) and actions by first projecting each input, assigning a type embedding to distinguish observations from actions, and interleaving the resulting embeddings to form the history representation $h_t$. The latent encoder $E_\theta$ encodes the local terrain height map using a CNN followed by spatial encoding, flattening, and projection; this representation is concatenated with the projected latest proprioceptive observation and passed through an MLP to produce the latent state embedding. Commands consist of the desired forward, lateral, and yaw velocities of the base \mbox{$g_t := (v^x_t, v^y_t, \omega_t)$}.  The future commands $g_{t:t+T-1}$ are concatenated to the inputs of the base policy $\pi_\theta$, reward model $R_\theta$, and value model $V_\theta$.












\begin{figure*}[t]
    \centering
    \includegraphics[width=.99\textwidth]{figs/model_qped.pdf}
    \caption{
    \footnotesize{Detailed world model architecture for the quadruped. 
    }
    }
     \vspace{-1em}
    \label{fig:model_qped}
\end{figure*}

\textbf{World Model Pretraining.}
We pretrain the world model for two epochs over the full simulation dataset $\simdataset$. Using a batch size of 512, this corresponds to approximately \num{3.69e5} gradient update steps. Optimization is performed with Adam using an initial learning rate of \num{2e-4}, which is annealed to \num{1e-4} via a cosine decay schedule, with a linear warmup over the first 10{,}000 steps. We apply data augmentation by injecting zero-mean Gaussian noise into the input observations (both proprioceptive and height map) during training. Pretraining requires approximately 28 hours on a single NVIDIA RTX 4500 Ada GPU.


\textbf{Hardware Deployment.} We use MPPI as implemented in TD-MPC \cite{hansen2022temporal} with the hyperparameters listed in \cref{tab:mppi_qped}. To encourage straight-line locomotion during quadruped experiments, we compute commanded base velocities $g_t$ from the robot’s current base pose using a PD controller on position, and provide these commands to the world model.


\textbf{Detailed Results.}
\Cref{tab:qpedallresults} reports detailed real-world quadruped locomotion results on both tasks across commanded forward speeds. We report both success rate (successful trials out of five) and average forward progress (mean $\pm$ standard deviation) over all trials at each speed. The \emph{Pretrained model} corresponds to zero-shot deployment of the simulation-pretrained world model without any real-world finetuning. While this model occasionally achieves partial forward progress, it fails to complete the task reliably, highlighting the severity of the sim-to-real dynamics gap. The \emph{Single-step BC policy}, which serves as the initial policy for the RLPD and IQL baselines prior to finetuning, improves performance in some settings but remains inconsistent and rarely achieves full task completion.

After real-world finetuning (32.1 minutes of data for Foam and 35.7 minutes for Slippery Slope), \Method\ consistently achieves the highest success rates and forward progress across all tested speeds. In contrast, both IQL and RLPD exhibit limited improvement despite access to the same real-world data budget. In particular, RLPD destabilized the robot during adaptation on the Foam task and is therefore not reported for that condition.







\subsection{Ablations Details}
\label{app:ablations}
Each ablation corresponds to a separately pretrained world model. Across all ablations, the same planning hyperparameters and evaluation environments are used, and each configuration is evaluated on an identical set of randomized environments. Next, we discuss platform specific details.

\textbf{Manipulation.} For both manipulation tasks, we evaluate success rates over the initial condition and domain randomization described in the environments from \cite{yin2026emergent}. 


\textbf{Quadruped.}
The robot is commanded to walk forward at a specified target speed until episode termination. For each evaluated model, we instantiate environments covering all combinations of parameters listed in \cref{tab:ablations_qped}, resulting in 1080 distinct environments per model. All models are evaluated on the same fixed set of environments to ensure fair comparison. Each environment is run for a single episode, during which state-based reward is accumulated until termination. Episodes terminate either after 1000 simulation steps or upon failure, defined as body contact with the ground or violation of base orientation limits. Results are summarized in \cref{tab:ablations}, where we report the average accumulated reward per episode across all environments.

